\documentclass{article}
\usepackage[T1]{fontenc}
\usepackage{lmodern}
\usepackage{graphicx}
\usepackage{amsmath,amssymb}
\usepackage[a4paper,margin=2cm]{geometry}
\usepackage[absolute,overlay]{textpos}
\setlength{\TPHorizModule}{1mm}
\setlength{\TPVertModule}{1mm}
\usepackage[indonesian]{babel}
\usepackage{hyperref}
\setlength{\emergencystretch}{2em}
% unit: Halaman 27

\begin{document}

% === Halaman 27 (llm) ===

\section*{Terminologi Steganografi}

\begin{enumerate}
  \item \textit{Embedded message (hiddentext)} \textit{atau secret message}: pesan yang disembunyikan.
  
  Bisa berupa teks, gambar, audio, video, dll
  
  \item \textit{Cover-object (covertext)}: pesan yang digunakan untuk menyembunyikan \textit{embedded message}.
  
  Bisa berupa teks, gambar, audio, video, dll
  
  \item \textit{Stego-object (stegotext)}: pesan yang sudah berisi pesan \textit{embedded message}.
  
  \item \textit{Stego-key}: kunci yang digunakan untuk menyisipan pesan dan mengekstraksi pesan dari stegotext.
\end{enumerate}

\end{document}
